{
\small
\renewcommand{\arraystretch}{1.15}
\setlist[enumerate]{leftmargin=*, nosep, labelsep=0.3em}

\begin{tabularx}{\textwidth}{|>{\raggedright\arraybackslash}p{3.5cm}|>{\raggedright\arraybackslash}X|}
\hline
\textbf{Tên Use case} & Duyệt bài thi \\
\hline
\textbf{ID} & UC-0010 \\
\hline
\textbf{Tác nhân} & Người dùng (công khai) \\
\hline
\textbf{Mô tả} & Người dùng duyệt và tìm kiếm danh sách bài thi đã xuất bản \\
\hline
\textbf{Điều kiện tiên quyết} & Không yêu cầu đăng nhập \\
\hline
\textbf{Điều kiện sau} & Danh sách bài thi được hiển thị cho người dùng \\
\hline
\textbf{Kích hoạt} & Người dùng mở trang danh sách bài thi \\
\hline
\textbf{Luồng chính} &
\begin{enumerate}
    \item Người dùng mở trang danh sách bài thi.
    \item Hệ thống hiển thị danh sách bài thi đã xuất bản với tên, thẻ phân loại và thống kê.
    \item Người dùng có thể lọc theo tên, thẻ phân loại.
    \item Người dùng có thể sắp xếp theo các tiêu chí khác nhau (tên, ngày tạo) với thứ tự tăng/giảm.
    \item Hệ thống hỗ trợ phân trang dạng cursor-based.
\end{enumerate} \\
\hline
\textbf{Luồng thay thế} &
\begin{enumerate}
    \item \textbf{Xem thống kê bài thi:}
    \begin{enumerate}
        \item[3A.] Người dùng chọn xem thống kê của một bài thi.
        \item[4A.] Hệ thống hiển thị: điểm trung bình, thời gian trung bình, tỷ lệ đúng theo từng thẻ phân loại.
    \end{enumerate}
    \item \textbf{Xem chi tiết câu hỏi:}
    \begin{enumerate}
        \item[3B.] Người dùng chọn xem chi tiết một câu hỏi.
        \item[4B.] Hệ thống hiển thị nội dung câu hỏi, các lựa chọn, giải thích và tệp đính kèm.
    \end{enumerate}
\end{enumerate} \\
\hline
\textbf{Luồng ngoại lệ} &
\begin{enumerate}
    \item \textbf{Không tìm thấy bài thi:}
    \begin{enumerate}
        \item[2Err.] Hệ thống trả về danh sách rỗng khi không có bài thi phù hợp với bộ lọc.
    \end{enumerate}
\end{enumerate} \\
\hline
\end{tabularx}
\captionof{table}{Đặc tả Use case: UC-0010 Duyệt bài thi}
}
